% RQ3 Discussion: Proposal Pipeline Effects

\subsection{Interpretation of Results}

\subsubsection{The Filtering Tradeoff}

Temp-checks function as noisy classifiers. Perfect filtering (passing only proposals that will succeed in formal voting) is impossible given:
\begin{itemize}
    \item Lower temp-check participation (noisier signal)
    \item Information revealed between temp-check and formal vote
    \item Different populations (temp-check vs. formal vote participants)
\end{itemize}

Our results show that moderate thresholds (20-30\%) provide good filtering while avoiding excessive false negatives. Thresholds above 40\% filter substantial numbers of ultimately viable proposals.

\subsubsection{Fast-Track as Acceleration}

Fast-tracks are most valuable for truly uncontroversial proposals where the pipeline stages add delay without adding information. Our 70\% threshold captures this use case.

However, fast-tracks create risks:
\begin{itemize}
    \item Capture opportunity: coordinate supporters for early vote, bypass deliberation
    \item Legitimacy questions: was there adequate consideration?
    \item Precedent: what counts as ``uncontroversial'' may expand over time
\end{itemize}

We recommend clear criteria for fast-track eligibility and periodic review of fast-tracked proposals' outcomes.

\subsubsection{Expiry as Forcing Function}

Expiry windows prevent indefinite proposal accumulation but may force premature resolution. Our results suggest:
\begin{itemize}
    \item 30 days is too aggressive for most governance contexts
    \item 60 days provides reasonable balance
    \item 90+ days may allow proposals to linger without resolution
\end{itemize}

The optimal window depends on proposal complexity. A tiered system (short expiry for simple, long for complex) may be valuable.

\subsection{Pipeline Design Principles}

Based on our results, we propose the following design principles:

\begin{enumerate}
    \item \textbf{Filter early, filter gently}: Use temp-checks with moderate thresholds to reduce formal voting load without excluding viable proposals

    \item \textbf{Accelerate consensus, not controversy}: Fast-track only proposals with clear supermajority support; never fast-track high-stakes decisions

    \item \textbf{Bound but don't compress}: Set expiry windows that prevent accumulation but allow adequate deliberation for the proposal type

    \item \textbf{Make stages meaningful}: Each pipeline stage should add information value; remove stages that are purely procedural delay
\end{enumerate}

\subsection{Comparison with Real Pipelines}

\subsubsection{Uniswap}

Uniswap's two-stage temp-check (temperature check + consensus check) aligns with our findings supporting moderate early filtering. However, the specific thresholds (majority for temperature, 67\% for consensus) may be calibrated differently than our recommendations.

\subsubsection{Compound}

Compound's single-stage process (no temp-check) maximizes inclusion but may result in more failed proposals consuming voter attention.

\subsubsection{Optimism}

Optimism's review periods function similarly to temp-checks but without explicit support thresholds. This provides filtering through community feedback rather than vote counts.

\subsection{Limitations}

\subsubsection{Simplified Proposal Quality}

We model proposals with fixed quality distributions. Real proposal quality is influenced by:
\begin{itemize}
    \item Proposer effort and expertise
    \item Community feedback during pipeline
    \item External events affecting relevance
\end{itemize}

\subsubsection{Homogeneous Stages}

Our model uses uniform stage parameters. Real pipelines often have stage-specific requirements (different quorums, durations) that we do not fully explore.

\subsubsection{No Learning}

Proposers do not learn from past failures. Real proposers may improve proposals based on feedback, affecting pipeline dynamics.

\subsection{Future Work}

\begin{enumerate}
    \item \textbf{Adaptive pipelines}: Simulate pipelines that adjust parameters based on proposal volume and type
    \item \textbf{Proposal quality dynamics}: Model how feedback improves proposal quality through pipeline stages
    \item \textbf{Proposer behavior}: Include strategic proposal timing and revision in response to pipeline incentives
    \item \textbf{Cross-DAO comparison}: Compare pipeline configurations across DAOs with different operational contexts
\end{enumerate}
